% Papiers de Sophie Germain — NAF 4073, lecture modernisée du volume entier.
%
% Pass: Opus 5 (claude-opus-5), 11 September 2026 — first pass, unchecked
% against the views by a human, and an interpretation rather than a record.
% Where this file and the transcriptions differ in substance, the
% transcriptions (batch-01.fr.tex … batch-03.fr.tex) are the record.
%
% Sources read: the three batch transcriptions of this volume and nothing
% else. No image was consulted, and no published edition — not Grun's, not
% Del Centina's, not the printed 1821 « Recherches ». All three batches were
% transcribed on Opus 5. The volume is complete: views 1–44, folios 1–32.
%
% What the volume turned out to be. An incoming dossier, and almost nothing
% else — the BnF's binding of letters addressed to Germain, out of Libri's
% papers. Of thirty-two folios, two are in her hand (the letter to Libri of
% 18 April 1831) and exactly one carries a formula. There is therefore no
% mathematical spine to follow, and this reading does not invent one: it
% follows the object. Three threads: the single identity on f. 32v; the
% mathematics her correspondents send her (Savart, Navier, Libri's heat
% memoir) without any of her own work being present; and the volume read as a
% dating instrument, the address leaves fixing what the undated billets do not
% say.
%
% Two datings are computed here and are this reading's own, not the leaves':
% the germinal note of Lagrange falls on 6 April 1796 or 7 April 1802 and on
% no other day of the Republican calendar; and the three Berlin leaves are
% three letters spanning 1776–1778, the last placed in July 1778 by Frederick
% II's entry into Bohemia. Both are footnoted at the point where they are made,
% with what would settle them.
%
% One correction to the transcriptions, argued rather than read: batch 2's
% note holds that « 19 = vendredi, 20 = samedi » corroborates the reading
% « 17 germinal (mercredi) ». That agreement is automatic and corroborates
% nothing about the numeral. The correction is made in the open, at § Lagrange.
%
% Notation translated here, each with a footnote at its first use: the leaf's
% (n²∓n)/2 into the triangular numbers T_{n−1}, T_n and its right-hand side
% into T_{n²}; Navier's and Germain's fourth-order equation of the elastic
% plate into the biharmonic Δ²; Fourier's heat equation set against a
% hyperbolic (relaxation) model to say where a derivation « from calorific
% vibrations » cannot be exact; « l'intégrale d'une équation différentielle
% linéaire » into the solution space and its dimension.
\documentclass[11pt,a4paper]{article}
% The preamble shared by every transcription of the mathematicians' manuscripts in Gallica.
%
% It rests on one idea: **uncertainty compounds, it does not resolve.** A
% transcription that smooths over a doubtful word has destroyed the only thing
% it was made for — a reader must be able to tell, at every point, what was on
% the page and what was supplied. The apparatus macros below are therefore the
% critical apparatus, not decoration.
%
% The same commands are recognised by `scripts/render.mjs`, which produces the
% site's reading view, and by `scripts/tei.mjs`. A macro added here must be
% added there too, or rendering fails loudly — which is the wanted behaviour.
%
% Comments here are English, like the rest of the repository. The documents
% this preamble sets are French, and that is deliberate: see the skills.

\ProvidesPackage{germain}[2026/09/15 Transcriptions of mathematicians' manuscripts in Gallica]

\usepackage{amsmath,amssymb,amsthm}
\usepackage{mathrsfs}
% Kept from the parent preamble so the renderer's diagram support stays
% available; these manuscripts draw no commutative diagrams, and nothing here uses it.
\usepackage{tikz-cd}
\usepackage[english,french]{babel}
\usepackage{fontspec}

% --- Greek ------------------------------------------------------------------
% The text face stays fontspec's default, Latin Modern, which has no Greek:
% XeTeX drops every glyph it lacks with a line in the log and nothing on the
% page, and the Greek words the manuscripts quote vanished from the PDFs. So
% the Greek and Greek Extended blocks get a XeTeX character class of their own,
% and entering or leaving a run of it switches the family to CMU Serif — the
% same Computer Modern design, with polytonic Greek — and back. Only the
% family changes: a Greek word in \textit or \textbf keeps its shape and series.
% The transitions to and from every other class carry the tokens that class
% has towards an ordinary letter, so French spacing before « ; » or inside
% « » still holds after a Greek word. No grouping, so no brace or \endgroup
% can come out unbalanced; maths is left alone, since XeTeX inserts nothing
% there. Kept identical in `exercices.sty`.
\makeatletter
\newfontfamily\germainGreekfont{cmun}[
  Extension=.otf, NFSSFamily=germaingreek,
  UprightFont=*rm, ItalicFont=*ti, SlantedFont=*sl,
  BoldFont=*bx, BoldItalicFont=*bi, BoldSlantedFont=*bl]
\newXeTeXintercharclass\germain@greek
\def\germain@greekrange#1#2{%
  \@tempcnta=#1\relax
  \loop
    \XeTeXcharclass\@tempcnta=\germain@greek
    \ifnum\@tempcnta<#2\advance\@tempcnta\@ne\repeat}
\germain@greekrange{"0370}{"03FF}
\germain@greekrange{"1F00}{"1FFF}
\def\germain@greekprev{\rmdefault}
\def\germain@greekin{%
  \edef\germain@greekprev{\f@family}\fontfamily{germaingreek}\selectfont}
\def\germain@greekout{\fontfamily{\germain@greekprev}\selectfont}
\def\germain@greekjoin{%
  \XeTeXinterchartokenstate=\@ne
  \@tempcnta=\z@
  \loop
    \ifnum\@tempcnta=\germain@greek\else
      \XeTeXinterchartoks\germain@greek\@tempcnta=\expandafter{%
        \expandafter\germain@greekout\the\XeTeXinterchartoks\z@\@tempcnta}%
      \XeTeXinterchartoks\@tempcnta\germain@greek=\expandafter{%
        \the\XeTeXinterchartoks\@tempcnta\z@\germain@greekin}%
    \fi
    \ifnum\@tempcnta<4095\advance\@tempcnta\@ne\repeat}
% babel-french sets its own classes, and saves and restores the interchar
% state, each time a language is selected: join again after it does.
\addto\extrasfrench{\germain@greekjoin}
\addto\extrasenglish{\germain@greekjoin}
\AtBeginDocument{\germain@greekjoin}
\makeatother

\usepackage{geometry}
\geometry{a4paper,margin=25mm,marginparwidth=28mm}
\usepackage{marginnote}
\usepackage{xcolor}
\usepackage[normalem]{ulem}
\setlength{\emergencystretch}{3em}
\usepackage{hyperref}
\hypersetup{colorlinks=true,linkcolor=black,urlcolor=black,pdfborder={0 0 0}}

% --- Batch metadata ---------------------------------------------------------
% Taken as they stand by the rendering script, which shows them at the head of
% the reading view. They fix what the file claims to cover. The macro names are
% the parent project's, so that the scripts stay shared; \folder here means the
% volume's slug — `fr-9115` — and \pages counts Gallica views.

\newcommand{\folder}[1]{\gdef\grothFolder{#1}}
\newcommand{\batch}[1]{\gdef\grothBatch{#1}}
\newcommand{\pages}[2]{\gdef\grothFirst{#1}\gdef\grothLast{#2}}
\newcommand{\foldertitle}[1]{\gdef\grothFoldertitle{#1}}

% The holder's own shelfmark, printed as they write it: « Français 9115 ».
\newcommand{\shelfmark}[1]{\gdef\grothShelfmark{#1}}

% The dating, copied verbatim from the holder's catalogue — never inferred
% here. « XIXe siècle » is the BnF's; a pass that dates a leaf from its content
% says so in a \note{}, as its own inference.
\newcommand{\dating}[1]{\gdef\grothDating{#1}}

% The status notice, printed in the title block and repeated as a diagonal
% watermark on every page of the PDF. The manuscripts are in the public
% domain; what this line declares is not a right but a quality — an
% unverified first pass by a machine. The skills write the line; a file
% without it gets no watermark, so leaving it out is visible in the output.
\newcommand{\watermark}[1]{\gdef\grothWatermark{#1}}

\gdef\grothFolder{?}\gdef\grothBatch{}\gdef\grothFirst{?}\gdef\grothLast{?}%
\gdef\grothFoldertitle{}\gdef\grothShelfmark{}\gdef\grothDating{}\gdef\grothWatermark{}

% --- The résumé -------------------------------------------------------------
\newenvironment{resume}
  {\small\setlength{\parskip}{0.45em}}
  {\par\medskip\hrule\medskip}

% --- Keywords ---------------------------------------------------------------
% In English, closing the modernised reading's résumé: the modern vocabulary
% under which to search for what the leaves construct. The manifest extracts
% them to tag the volume — this line is the single source of the tags.
\newcommand{\keywords}[1]{\par\smallskip\noindent
  {\footnotesize\textsc{Keywords} --- \textit{#1}}\par}

% --- The critical apparatus -------------------------------------------------

% The Gallica view number, in the margin. This is the anchor that lets a
% disagreement over a formula be settled by looking at *one* image rather than
% a volume of 750. The site uses it to turn the facsimile as the transcription
% is scrolled. A view is one scanned image: usually one side of a leaf.
\newcommand{\page}[1]{%
  \marginnote{\footnotesize\color{gray!70}\textsf{v.\,#1}}%
  \label{page:#1}\ignorespaces}

% The BnF's pencilled foliation, where the leaf carries it and a pass can read
% it: « 198r ». This is what the literature cites — « ff. 198r–208v » — and
% the one line that turns a citation into a view. Recorded, never inferred:
% a folio a pass cannot read is not written.
\newcommand{\folio}[1]{%
  \marginnote{\footnotesize\color{gray!50}\textsf{f.\,#1}}\ignorespaces}

% The modernised reading's coarser counterpart to \page{}: which views a
% section is drawn from.
\newcommand{\pagerange}[2]{%
  \marginnote{\footnotesize\color{gray!70}\textsf{v.\,#1--#2}}%
  \ignorespaces}

% Illegible: never guessed.
\newcommand{\ill}{\textcolor{red!60!black}{[\,\ldots\,]}}

% A doubtful reading: the word is given, and flagged as uncertain.
\newcommand{\uncertain}[1]{\uline{#1}}

% An editorial addition — a missing letter, an implied word.
\newcommand{\add}[1]{\textcolor{blue!50!black}{[#1]}}

% Struck out by the writer. What was crossed out is part of the document.
\newcommand{\struck}[1]{\sout{#1}}

% The transcriber's note, on the state of the page or the mathematics.
\newcommand{\note}[1]{\par\smallskip{\small\color{gray!60!black}\textit{[#1]}}\par\smallskip}

% A marginal note of hers, distinct from the body of the text.
\newcommand{\marginal}[1]{\par\smallskip\noindent\hspace{1em}%
  {\small\color{gray!60!black}#1}\par\smallskip}

% --- Title ------------------------------------------------------------------
% The title block is French, like the documents themselves.

\AddToHook{shipout/background}{%
  \ifx\grothWatermark\empty\else
    \begin{tikzpicture}[remember picture, overlay]
      \node[rotate=52, scale=3.4, align=center, text=gray!18,
            font=\sffamily\bfseries]
        at (current page.center) {\grothWatermark};
    \end{tikzpicture}%
  \fi}

\AtBeginDocument{%
  \begin{center}
    {\large\bfseries Manuscrits de mathématiciens (Gallica) ---
      \ifx\grothShelfmark\empty\grothFolder\else\grothShelfmark\fi}\\[2pt]
    {\itshape\grothFoldertitle}\\[4pt]
    {\small vues \grothFirst--\grothLast
      \ifx\grothBatch\empty\else\ (lot \grothBatch)\fi}\\[2pt]
    \ifx\grothDating\empty\else
      {\small\color{gray!60!black}Datation du catalogue : \grothDating}\\[2pt]
    \fi
    \ifx\grothWatermark\empty\else
      {\small\bfseries\color{red!45!gray}\grothWatermark}\\[2pt]
    \fi
    {\footnotesize\color{gray!60!black}Transcription établie d'après les images IIIF de Gallica.
     Source gallica.bnf.fr / Bibliothèque nationale de France.\\
     Les lectures incertaines sont soulignées, les illisibles notées [\,\ldots\,],
     les ajouts entre crochets ; \textsf{v.} numérote les vues de Gallica,
     \textsf{f.} la foliotation au crayon de la BnF.}
  \end{center}
  \vspace{1em}}

\endinput


\folder{naf-4073}
\shelfmark{NAF 4073}
\pages{1}{44}
\dating{XVIIIe siècle}
\watermark{Lecture automatique\\interprétation non vérifiée}
\foldertitle{Correspondance de Sophie GERMAIN — lecture modernisée du volume entier}

\begin{document}

\section*{Résumé}

\begin{resume}
Ce volume est un dossier entrant. Quarante-quatre vues, trente-deux feuillets,
reliés par la Bibliothèque à partir des papiers de Libri : ce sont les lettres
qu'elle a \emph{reçues} — onze billets de Fourier, deux de Delambre, un de
Lalande, un de Libri, un de Choron, un de Lagrange —, trois feuillets de
Lagrange écrits de Berlin en 1776--1778 qui ne lui sont pas adressés du tout,
une lettre anonyme de 1814, et quatre pages de sa main, la lettre à Libri du
18 avril 1831. Un seul feuillet de tout le volume porte une formule.

Il faut donc dire d'emblée ce que cette lecture ne peut pas faire. Il n'y a
pas ici de fil mathématique à suivre, parce qu'il n'y a pas de mathématiques à
suivre : ni le mémoire de 1811, ni celui de 1813, ni les \emph{Recherches} de
1821, ni une ligne de son arithmétique. Ce volume ne dit pas ce qu'elle
cherchait ; il dit ce qu'on lui envoyait, ce qu'on la pressait de publier, et
ce qu'elle pensait, à la fin, de l'état de la science.

La seule mathématique du volume tient en une ligne, écrite tête-bêche au haut
d'un feuillet d'adresse, par une main qui n'est pas celle de la lettre, après
le 13 mars 1814 : « le triangulaire d'un quarré est la somme des quarrés des
deux Triangulaires consécutifs », avec la formule qui la porte. En écrivant
$T_k = k(k+1)/2$, c'est l'identité $T_{n-1}^2 + T_n^2 = T_{n^2}$. Elle est
vraie, et elle se démontre en une ligne dès qu'on remarque les deux faits que
la phrase française contient déjà : les deux triangulaires consécutifs ont
pour somme $n^2$ et pour différence $n$, et l'identité du parallélogramme
$a^2+b^2 = \tfrac12((a+b)^2+(b-a)^2)$ fait le reste. On ignore de qui est la
main ; la transcription note sa ressemblance avec celle de 1831 et refuse de
conclure sur quinze mots, et cette lecture ne conclut pas davantage.

Autour de cette ligne, les lettres portent des mathématiques qui sont celles
des autres. Fourier lui fait passer le mémoire de Savart sur les vibrations
des lames élastiques et celui, lithographié, de Navier sur la flexion des
plans élastiques — c'est-à-dire les deux moitiés, l'expérience et l'analyse,
de la question qui l'occupait ; le nom moderne de leur objet commun est
l'opérateur biharmonique, et l'équation du mouvement libre d'une plaque mince
est $\partial^2 z/\partial t^2 + k^2\,\Delta^2 z = 0$. Fourier, en septembre
1820, l'invite « instamment à publier ». Delambre, en juillet 1821, accuse
réception d'un mémoire où elle combat les principes d'un « Géomètre célèbre »
qu'il ne nomme pas, et avoue que personne n'est en état de juger. Libri, de
Florence en 1825, annonce deux choses dont l'une est fausse — qu'on peut
donner à l'intégrale d'une équation différentielle linéaire un nombre
quelconque de constantes arbitraires, alors que leur nombre est exactement
l'ordre de l'équation — et dont l'autre, une dérivation des équations de la
chaleur à partir de vibrations, se heurte à ce qu'on formule aujourd'hui comme
la différence entre une équation parabolique et une équation hyperbolique.

Le volume rend aussi un service qu'on n'attend pas d'une correspondance : il
date. Presque chaque lettre a gardé sa couverture, et les adresses successives
— rue Sainte-Croix-de-la-Bretonnerie jusqu'en 1814, rue de Braque à partir de
1820 au plus tard — situent les onze pièces qui ne portent qu'un jour de la
semaine. Deux dates se calculent : le billet de Lagrange « ce 17 germinal
(mercredi) » ne peut être que du 6 avril 1796 ou du 7 avril 1802, ces deux
jours étant les seuls mercredis 17 germinal de tout le calendrier
républicain ; et les trois feuillets de Berlin, que la reliure donne à la
file, sont trois lettres distinctes dont la dernière est de juillet 1778, date
à laquelle Frédéric II entre en Bohême.

De ce qu'elle pensait, elle a laissé ici son propre compte rendu, et il tient
en une phrase écrite dix semaines avant sa mort : « décidément il y a un tort
sur tout ce qui tient aux mathématiques votre préoccupation celle de Cauchy,
la mort de Mr fourier, pour […] cet élève Gallois qui malgré son impertinence
annonçoit des dispositions heureuses en a tant fait qu'il a été chassé de
l'école normale ». Libri en exil, Cauchy parti, Fourier mort, Galois chassé :
elle fait le compte, et elle range le dernier avec les autres.

Où cela mène, et sous quels noms chercher. Du côté de la ligne du feuillet
32v : les nombres figurés, l'identité de Nicomaque $1^3+\cdots+n^3 = T_n^2$,
dont on tire ici le corollaire $T_{n^2} = 2T_n^2 - n^3$, et plus largement les
identités entre nombres triangulaires et sommes de deux carrés. Du côté des
plaques : l'opérateur biharmonique, la théorie des plaques minces de
Kirchhoff et Love, les figures de Chladni et les lignes nodales comme lieux
des zéros des modes propres. Du côté de la lettre de Libri : la dimension de
l'espace des solutions d'une équation linéaire, et les modèles hyperboliques
de la conduction de la chaleur. Et pour la lettre du 18 avril 1831, une
histoire qui n'est pas mathématique et que les mathématiciens lisent quand
même.

\keywords{triangular numbers, figurate number identities, sum of two squares,
Nicomachus's theorem, biharmonic operator, Kirchhoff–Love plate theory,
Chladni figures, nodal lines, heat equation, hyperbolic heat conduction,
linear ODE solution space}
\end{resume}

\section*{Ce que le volume contient}

\pagerange{1}{44}
Quarante-quatre vues, trente-deux feuillets. Les quatre premières vues et les
deux dernières sont la reliure, les gardes et les fiches de la
Bibliothèque\footnote{Vue 1, plat supérieur ; vue 2, contre-plat et l'étiquette
de cote ; vue 3, les fiches collées sur la garde, « Volume de 32 feuillets,
15 Mars 1875 » ; vue 4, une table du contenu au crayon d'une main de
catalogueur ; vues 43 et 44, gardes blanches et plat inférieur. Rien de tout
cela n'est transcrit, et rien n'en est de la main d'aucun des correspondants.} ;
entre elles, trente-deux feuillets montés un à un, la plupart des vues
photographiant deux feuillets à la fois. Le volume vient des papiers de Libri.

Ce n'est pas un cahier de mathématiques, et ce n'est pas non plus, comme
Français 9118, un dossier mêlé où l'on trouverait ses brouillons à côté des
lettres reçues. C'est un dossier \emph{entrant}, et presque rien d'autre : sur
les trente-deux feuillets, deux seulement portent son écriture — les quatre
pages de la lettre à Libri du 18 avril 1831 —, et un seul feuillet de tout le
volume porte une formule.

Les pièces se rangent ainsi.

\begin{itemize}
  \item \emph{Onze billets et lettres de Fourier}, dont deux datés (29
    septembre 1820, 30 janvier 1827) et les autres portant un jour de la
    semaine et rien de plus : « vendredi matin », « mardi soir », « jeudi …
    mai », « mardi matin », « lundi matin ». Ils occupent, avec leurs
    couvertures, les vues 12 à 27.
  \item \emph{Deux lettres de Delambre} : 14 mai 1810, sur la médaille de
    Lalande décernée à Gauss, et 21 juillet 1821, sur un mémoire d'elle et sur
    les places aux séances publiques de l'Académie.
  \item \emph{Une lettre de Libri}, Florence, 17 novembre 1825.
  \item \emph{Une lettre d'elle à Libri}, Paris, 18 avril 1831, signée.
  \item \emph{Un billet de Lagrange} daté « 17 germinal », et \emph{trois
    feuillets de Lagrange écrits de Berlin} qui ne lui sont pas adressés et
    qui sont antérieurs de vingt ans à tout le reste.
  \item \emph{Un billet de Lalande}, 4 novembre 1797 : la pièce la plus
    ancienne du volume qui la concerne, et des excuses.
  \item \emph{Un billet de Choron}, 29 janvier 1818, sur l'enseignement
    simultané de la lecture et de l'écriture en musique.
  \item \emph{Une lettre non signée} du 13 mars 1814, en deux feuillets,
    d'une main qui n'est celle d'aucun des autres, et dont la couverture porte
    la seule formule du volume.
\end{itemize}

Deux de ces correspondants ne figurent pas dans la notice de la
Bibliothèque, qui énumère Delambre, Fourier, Germain, Lagrange, Lalande et
Libri : Choron, et l'inconnu de 1814.\footnote{Le constat est celui de la
transcription, au feuillet 1r ; il est répété ici parce qu'il porte sur
l'objet et non sur le texte. La datation du catalogue, « XVIIIe siècle »,
copiée telle quelle en tête de cette lecture, ne convient qu'aux feuillets de
Lagrange (1776--1778) et au billet de Lalande (1797) : tout le reste du volume
est du XIXe.}

Presque chaque lettre a gardé son feuillet d'adresse, plié et cacheté. Ces
couvertures, qui ne disent rien de mathématique, sont l'instrument le plus
utile du volume : elles datent ce qui n'est pas daté. On y revient plus bas.

\section*{La seule mathématique du volume : le triangulaire d'un carré}

\pagerange{42}{42}
Au haut du feuillet d'adresse de la lettre non signée du 13 mars 1814, écrite
tête-bêche par rapport à l'adresse et d'une main qui n'est pas celle de la
lettre, une seule ligne de mathématiques, avec la phrase qui la lit :
« le triangulaire d'un quarré est la somme des quarrés des deux Triangulaires
consécutifs ».\footnote{Le dernier mot est pris dans le fond du cahier ; la
transcription lit « consécu » et supplée le reste. C'est le seul mot supplé
de toute la note.}

\subsection*{L'énoncé}

Le feuillet écrit
\[
\left(\frac{n^2-n}{2}\right)^{\!2} + \left(\frac{n^2+n}{2}\right)^{\!2}
= \frac{n^4+n^2}{2}.
\]
En notant $T_k = k(k+1)/2$ le $k$-ième nombre triangulaire, les deux quantités
élevées au carré sont $T_{n-1}$ et $T_n$, et le second membre est $T_{n^2}$.
L'énoncé est donc
\[
T_{n-1}^{\,2} + T_n^{\,2} = T_{n^2},
\]
c'est-à-dire exactement ce que dit la phrase française : le triangulaire d'un
carré est la somme des carrés des deux triangulaires consécutifs.\footnote{La
phrase est plus précise que la formule ne le laisse voir. « Les deux
Triangulaires consécutifs » ne sont pas deux triangulaires quelconques qui se
suivent : ce sont les deux dont la somme est le carré en question, puisque
$T_{n-1} + T_n = n^2$. La formule fixe lesquels, la phrase ne le fait pas.}
Sous forme binomiale,
$\binom{n}{2}^{\!2} + \binom{n+1}{2}^{\!2} = \binom{n^2+1}{2}$.

\subsection*{Elle est vraie, et d'une ligne}

Posons $a = T_{n-1}$ et $b = T_n$. Deux faits immédiats :
\[
a + b = \frac{n^2-n}{2} + \frac{n^2+n}{2} = n^2,
\qquad
b - a = n.
\]
L'identité du parallélogramme --- $a^2 + b^2 = \tfrac12\big((a+b)^2 +
(b-a)^2\big)$, qui n'est que le développement des deux carrés --- donne alors
\[
a^2 + b^2 = \tfrac12\big(n^4 + n^2\big) = \frac{n^2(n^2+1)}{2} = T_{n^2}.
\]
C'est tout. La démonstration ne demande rien qui ne fût disponible en 1814, et
la disposition même de la phrase française --- un carré, ses deux triangulaires
--- est celle du calcul : le carré est la \emph{somme} des deux triangulaires,
leur \emph{différence} est sa racine, et l'identité du parallélogramme fait le
reste.

Une conséquence, qu'aucune main ne tire ici mais qui vaut d'être écrite.
Nicomaque donne $1^3 + 2^3 + \cdots + n^3 = T_n^{\,2}$ ; en portant cela dans
l'identité, $T_{n-1}^{\,2} = T_n^{\,2} - n^3$, donc
\[
T_{n^2} = 2\,T_n^{\,2} - n^3 .
\]
Le triangulaire d'un carré vaut deux fois la somme des $n$ premiers cubes,
moins le dernier.

\subsection*{La ligne de trois nombres}

Sous l'égalité, la note porte une ligne isolée, trois nombres reliés par un
trait horizontal continu, que la transcription lit $1 + 64 + 49$ en avertissant
que le second signe est empâté.\footnote{La transcription ajoute : « Rien ne
rattache cette ligne à l'égalité qui la précède. » La présente lecture n'en
rattache pas davantage, et ce qui suit est un constat, non une
reconstruction.}

Aucune lecture de ces trois nombres n'instancie l'identité. La somme vaut
$114$, qui n'est ni triangulaire --- $T_{14} = 105$, $T_{15} = 120$ --- ni
carré. Les plus petits cas de l'identité, ceux qu'on écrirait pour la vérifier,
sont $1 + 9 = 10$ pour $n = 2$ et $9 + 36 = 45$ pour $n = 3$, et ce n'est ni
l'un ni l'autre. Les trois nombres sont les carrés de $1$, $7$ et $8$ ; si le
second signe est un signe d'égalité plutôt qu'un plus, $1 + 64 = 49$ est faux,
en sorte que la lecture de la transcription est la seule qui fasse de la ligne
un énoncé arithmétique recevable --- une addition posée. Ce qu'elle additionne,
le feuillet ne le dit pas.

\subsection*{La main, et la date}

Deux choses sont sûres, une troisième ne l'est pas. Sûres : la note a été
écrite sur ce feuillet \emph{après} qu'il eut servi de couverture, donc après
le 13 mars 1814 ; et elle n'est pas de la main qui a écrit la lettre, puisque
l'encre et le module diffèrent et que l'écriture est tête-bêche par rapport à
l'adresse. Non sûre : de qui elle est. La transcription note que la main
ressemble à celle de la lettre du 18 avril 1831 --- même inclinaison, même
$d$ bouclé, mêmes longues liaisons plates --- et refuse d'en conclure, quinze
mots n'étant pas une attribution. Cette lecture ne conclut pas
davantage.\footnote{Ce qui trancherait : une comparaison des quinze mots avec
un corpus daté de sa main --- les brouillons de Français 9118, par exemple ---
sur les lettres que la note contient en propre, le $q$ de « quarré », le $T$
majuscule de « Triangulaires », le $C$ de « consécutifs ». La présente
lecture ne dispose que des transcriptions, et une transcription ne conserve pas
la forme des lettres.}

\section*{Les surfaces élastiques : ce qu'on lui envoie, et ce qu'on la presse de publier}

\pagerange{10}{19}
Le volume ne contient pas une ligne de sa théorie des plaques. Il contient ce
qui l'entoure : deux correspondants qui la poussent à publier, et un troisième
qui lui fait porter les mémoires des autres.

\subsection*{Fourier, 29 septembre 1820}

Fourier part quinze jours à la campagne et écrit avant de partir :
\begin{quote}
Je suis persuadé que vous aurez donné de nouveaux développemens à la théorie
dont vous vous occupez, et que je vous invite instamment à publier. personne
ne peut traiter avec plus de succès cette difficile et intéressante question.
\end{quote}
« La théorie dont vous vous occupez » n'est pas nommée. Neuf mois plus tard,
Delambre accuse réception d'un mémoire ; les \emph{Recherches sur la théorie
des surfaces élastiques} paraissent en 1821. Le rapprochement est
vraisemblable et il n'est pas écrit sur le feuillet : ni Fourier ni Delambre ne
nomment l'ouvrage, et cette lecture ne le nomme qu'en le donnant pour ce qu'il
est, une conjecture.\footnote{C'est aussi ce que fait la transcription, qui
écrit en note, au feuillet 4r : « le mémoire est vraisemblablement les
“Recherches sur la théorie des surfaces élastiques”, parues en 1821 ; le
rapprochement est celui de la présente lecture et non du document ». Les deux
lectures, indépendamment, s'arrêtent au même point.}

\subsection*{Delambre, 21 juillet 1821, et le « Géomètre célèbre »}

\begin{quote}
il ne m'appartient pas et peut-être n'appartient-il à personne de prononcer sur
les questions que vous regardez comme douteuses. mais ce qui obtiendra une
approbation générale c'est la réserve avec laquelle vous exposez vos opinions
et vos conjectures, ce sont les égards avec lesquels vous traitez le Géomètre
célèbre dont vous combattez les principes qui vous ont paru ne pouvoir
s'accorder avec ce que vous démontrez.
\end{quote}
Trois choses se lisent là. D'abord qu'elle a annoncé elle-même ses conclusions
comme douteuses. Ensuite que le secrétaire perpétuel se déclare hors d'état de
juger --- « peut-être n'appartient-il à personne » --- ce qui, en 1821 et sur
la question des plaques, est un aveu et non une politesse. Enfin qu'elle
combat les \emph{principes} de quelqu'un, c'est-à-dire son point de départ
physique, et non un résultat.

Qui est ce géomètre ? Le feuillet ne le nomme pas. Le candidat que la
disposition de la question rend le plus naturel est Poisson, dont le mémoire de
1814 tire l'équation des surfaces élastiques d'une hypothèse moléculaire, et
dont les principes sont précisément ce qu'un adversaire attaquerait avant les
conséquences ; mais c'est une conjecture de cette lecture, elle n'est pas dans
le feuillet, et elle ne s'appuie sur aucune pièce du volume.\footnote{Ce qui
trancherait est hors du volume : le texte du mémoire que Delambre accuse
réception, qui n'y est pas. Le rapprochement est donné ici parce qu'un lecteur
le fera de toute façon, et qu'il vaut mieux qu'il le lise annoncé comme
conjecture.}

\subsection*{Savart, Navier, et l'équation}

Un billet de Fourier sans date, « mardi soir », fait passer deux
imprimés\,:\footnote{Sur la datation de ce billet, voir plus bas : les
feuillets d'adresse et la santé de sa mère le placent entre 1820 et 1823.}
\begin{quote}
j'ai l'honneur d'adresser à mademoiselle germain le dernier numéro des annales
de physique où se trouve le mémoire de Mr Savart sur les vibrations des lames
élastiques […] j'espère lui envoyer demain un mémoire qui a été présenté à
l'académie par Mr Navier et qui a pour objet l'analyse des flexions des plans
élastiques. l'auteur a fait copier ce manuscrit par les presses
lithographiques.
\end{quote}
Les deux envois sont les deux moitiés de la question telle qu'elle se posait
alors : l'expérience et l'analyse. Savart continue les figures de Chladni ---
le sable qui se rassemble sur les lignes où la plaque ne bouge pas ; Navier
donne l'équation de la flexion.

En termes d'aujourd'hui, l'objet commun est l'opérateur bi\-harmonique. Pour une
plaque mince, homogène et isotrope, d'épaisseur $h$ et de déplacement
transversal $w(x,y)$, la flexion statique sous une charge $q$ est régie par
\[
D\,\Delta^2 w = q,
\qquad
D = \frac{E h^3}{12\,(1-\nu^2)},
\]
et le mouvement libre par
\[
\frac{\partial^2 z}{\partial t^2} + k^2\,\Delta^2 z = 0 ,
\]
$\Delta^2 = \Delta\Delta$ étant le bilaplacien.\footnote{Les noms sont
postérieurs aux feuillets : « plaque de Kirchhoff--Love » renvoie à Kirchhoff
(1850) et à Love (1888), et le module de rigidité $D$ sous la forme écrite ici,
avec le coefficient de Poisson $\nu$, appartient à la seconde moitié du siècle.
Ni Navier ni Germain n'écrivent $\Delta^2$ ; ce que Navier écrit est une
équation du quatrième ordre en coordonnées cartésiennes, et la forme compacte
est une commodité de notation moderne, non une découverte qu'on leur prête.}
Les lignes nodales que le sable de Savart dessine sont les lieux $z = 0$ des
modes propres, c'est-à-dire les zéros des solutions de $\Delta^2 \phi =
\lambda \phi$ sous les conditions de bord de la plaque --- lesquelles ont
occupé le reste du siècle, et sont la vraie difficulté que ni ce billet ni ce
volume n'abordent.

Il faut dire ce que le volume ne permet pas de dire. Aucune des pièces
mathématiques de Germain sur les plaques n'est ici : ni le mémoire de 1811, ni
celui de 1813, ni les \emph{Recherches} de 1821. La question de savoir où son
analyse a manqué ne se pose donc pas sur ces feuillets, et cette lecture ne la
pose pas.\footnote{Elle se pose ailleurs, et le volume Français 9118 en porte
la trace. Le renvoi est fait ici pour marquer la limite de ce que NAF 4073
établit, non pour importer une conclusion : rien de ce qui suit dans cette
lecture ne s'appuie sur un feuillet qui n'est pas dans ce volume.}

\section*{Libri : la chaleur, les constantes arbitraires, et le rapport que Cauchy doit depuis six mois}

\pagerange{16}{39}
La lettre de Florence du 17 novembre 1825 est la seule pièce du volume, en
dehors de la note du feuillet 32v, qui avance un énoncé mathématique.

\subsection*{« Un nombre quelconque de constantes arbitraires »}

\begin{quote}
J'ai envoyé un mémoire à l'Institut sur la théorie de la chaleur, ou je crois
avoir démontré qu'on peut donner à l'Intégrale d'une équation différentielle
linéaire un nombre quelconque de constantes arbitraires.
\end{quote}
Tel quel, l'énoncé est faux, et il faut le dire. Pour une équation
différentielle linéaire d'ordre $n$ à coefficients continus sur un intervalle,
l'ensemble des solutions est un espace vectoriel de dimension exactement $n$ :
c'est le théorème de Cauchy--Lipschitz appliqué au système du premier ordre
équivalent, l'application qui à une solution associe le vecteur
$\big(y(t_0), y'(t_0), \dots, y^{(n-1)}(t_0)\big)$ étant un
isomorphisme.\footnote{Le nom de Cauchy--Lipschitz est ici anachronique de
peu : le théorème d'existence et d'unicité sous condition de Lipschitz est de
Cauchy, dans les années 1820, et Lipschitz l'affine en 1876. Pour les seules
équations linéaires, le fait que l'intégrale générale dépende de $n$
constantes et de $n$ seulement était su et admis en 1825, ce qui est
précisément ce qui rend l'annonce de Libri remarquable.} Le nombre de
constantes arbitraires \emph{essentielles} est donc un invariant de
l'équation, égal à son ordre, et aucune démonstration ne peut l'augmenter.

Ce qui peut être vrai, et qui est la seule façon de sauver la phrase, est
autre chose : une \emph{représentation} d'une solution peut contenir autant de
paramètres qu'on veut sans que l'ensemble décrit s'agrandisse --- des
constantes non essentielles, dont la variation ne change pas la fonction,
comme il en apparaît naturellement dans les solutions écrites sous forme
d'intégrale définie, où le contour et les bornes sont à disposition. Que ce
soit là ce que Libri démontre, le feuillet ne le dit pas : il donne l'annonce
et non la preuve, le mémoire n'est pas dans le volume, et cette lecture ne le
juge pas sur une phrase de lettre.

\subsection*{« L'hypothèse des vibrations calorifiques »}

\begin{quote}
J'ai aussi déduit les équations du mouvement de la chaleur de l'hypothèse des
vibrations calorifiques.
\end{quote}
C'est le programme qui oppose, dans les années 1820, une chaleur conçue comme
vibration de la matière à la chaleur-fluide dont Fourier avait su se passer.
Le point de difficulté, tel qu'on le formule aujourd'hui, est que les deux
hypothèses ne conduisent pas au même \emph{type} d'équation. L'équation de
Fourier
\[
\frac{\partial u}{\partial t} = k\,\Delta u
\]
est parabolique : irréversible, et de vitesse de propagation infinie. Une
vibration donne une équation hyperbolique, réversible et de vitesse finie. On
ne passe de la seconde à la première par aucun calcul exact, mais seulement par
un passage à la limite : en écrivant, avec un temps de relaxation $\tau$,
\[
\tau\,\frac{\partial^2 u}{\partial t^2} + \frac{\partial u}{\partial t}
= k\,\Delta u ,
\]
on retrouve l'équation de Fourier quand $\tau \to 0$, et l'équation des ondes
quand la conduction est négligeable.\footnote{Cette équation est celle de
Cattaneo (1948) et Vernotte (1958) ; elle est citée ici pour dire précisément
où gît la difficulté que la lettre de 1825 ne mentionne pas, et non pour
suggérer que Libri en disposait. Il ne pouvait pas en disposer.} Ce que la
lettre annonce est donc une dérivation dont on sait aujourd'hui qu'elle ne peut
pas être exacte ; ce qu'elle valait dans le mémoire, le volume ne permet pas
de le dire.

\subsection*{Le rapport de Cauchy, et les billets de Fourier}

La lettre se termine par une demande précise :
\begin{quote}
Voudriez vous avoir la bonté de chatouiller un peu Mr. Fourier afin qu'il fasse
hâter le rapport que Mr. Cauchy doit faire depuis six mois sur le mémoire sur
la théorie des Nombres que vous avez lu ?
\end{quote}
Deux faits de fait, qui ne demandent aucune interprétation : elle a lu le
mémoire de Libri sur la théorie des nombres, et Libri la tient pour capable de
faire bouger le secrétaire perpétuel de l'Académie.

Deux billets non datés de Fourier semblent répondre à cette démarche. Dans
l'un, « lundi matin » : « je n'ai point trouvé chez moi le mémoire de mr.
libri. […] je vais aujourd'hui à l'académie et je me ferai remettre cet ouvrage
que j'enverrai à mademoiselle germain. » Dans l'autre, « vendredi matin » : « je
m'acquitterai de la recommandation que vous me confiez —
me mande Mr Libri. il m'en avoit aussi parlé dans sa dernière
lettre. je ferai mon possible pour que ce vœu ne soit pas différé. il nous a
envoyé un nouveau mémoire sur la théorie de la chaleur. »\footnote{Les deux
mots soulignés sont donnés par la transcription comme d'une lecture douteuse,
et la construction de la phrase n'a pas pu y être rétablie. Le sens général ---
une recommandation transmise par elle, que Libri a lui aussi demandée par
lettre --- ne dépend d'aucun des deux mots.} Que ces billets soient la suite de
la lettre du 17 novembre est une conjecture de cette lecture : elle repose sur
la concordance de trois éléments --- un mémoire de Libri sur la chaleur à
l'Institut, une démarche qu'elle transmet, un vœu qu'on ne veut pas voir
différé --- et sur rien qui soit écrit.

\section*{Gauss, la médaille de Lalande et la pendule}

\pagerange{7}{9}
Le 14 mai 1810, Delambre lui écrit d'une chose qui n'est pas des mathématiques
et qui en dit long. Gauss vient de recevoir la médaille fondée par Lalande ;
sur les 500 francs qu'elle vaut, 380 restent en argent, et Gauss préférerait
une pendule. Delambre cite ses mots :
\begin{quote}
qu'il soit de 60 ou de 300 fr. cela m'est indifférent pourvû que la montre soit
assez élégante pour être offerte en cadeau à mon épouse […] peut-être
Mademoiselle Sophie Germain ( à laquelle je vous prie de faire mille complimens
de ma part ) auroit la bonté de se charger du choix.
\end{quote}
Delambre corrige au passage l'allemand de Gauss : « il a traduit par
\emph{montre à Pendule} l'expression allemande Pendeluhr », et ce que Gauss
veut est une pendule.

Ce qui compte ici est la date. En mai 1810, Gauss demande qu'on s'adresse à
elle pour un cadeau destiné à sa femme : il sait depuis quelques années qui
était « M. Le Blanc », et la lettre le montre traitant avec elle sur le pied de
la familiarité et non de la déférence savante.\footnote{La date à laquelle
Gauss a appris son nom n'est pas dans ce volume, et cette lecture ne l'y fait
pas entrer : elle relève de l'épisode Pernety, pendant l'occupation de
Brunswick, dont les traces sont ailleurs. Ce que NAF 4073 établit, tout seul,
est qu'en mai 1810 Gauss la nomme spontanément à Delambre comme la personne à
qui confier ce choix.}

\section*{Lagrange : un billet de germinal, et trois feuillets de Berlin}

\pagerange{29}{34}

\subsection*{Le billet de germinal, et sa date}

\begin{quote}
Paris ce 17 germinal (mercredi)

Lagrange présente ses respects à Mademoiselle Germaine. […] il se fait un
devoir de la prévenir qu'il sera à ses ordres le 19 et le 20. (vendredi, et
samedi)
\end{quote}
Le billet donne un quantième républicain et trois jours de la semaine. Cela
suffit presque à le dater, et il vaut de le faire proprement, parce que la
manière naturelle de le faire est fausse.

La transcription note que « la lecture “17” s'accorde avec le reste du billet,
où le 19 est un vendredi et le 20 un samedi ». Cet accord est automatique : si
le 17 est un mercredi, le 19 est un vendredi et le 20 un samedi, quelle que
soit l'année et quel que soit le mois. Les trois indications n'en font qu'une,
et elle ne confirme pas le chiffre 17 --- elle confirme seulement que les trois
quantièmes ont été lus de façon cohérente entre eux.\footnote{C'est le seul
point où cette lecture corrige la transcription, et elle le corrige sur un
raisonnement, non sur une lecture : le feuillet dit ce que la transcription dit
qu'il dit. La remarque vaut d'être faite parce que l'argument d'accord, s'il
était laissé tel quel, donnerait à la lecture « 17 » un appui qu'elle n'a pas.}

La contrainte utile est donc unique : \emph{le 17 germinal tombe un mercredi}.
Le calendrier républicain n'a que quatorze années, et sur ces quatorze, deux
seulement la satisfont :
\begin{itemize}
  \item 17 germinal an IV $=$ mercredi 6 avril 1796 ;
  \item 17 germinal an X $=$ mercredi 7 avril 1802.
\end{itemize}
Rien dans le billet ne sépare les deux. Elle a vingt ans dans le premier cas,
vingt-six dans le second ; Lagrange est à Paris dans les deux. Ce que cette
lecture peut affirmer est donc : le billet est du 6 avril 1796 ou du 7 avril
1802, et d'aucun autre jour.\footnote{Le calcul est celui de cette lecture. Il
suppose la concordance usuelle du calendrier républicain (1\textsuperscript{er}
vendémiaire de l'an IV $=$ 23 septembre 1795, de l'an X $=$ 23 septembre 1801),
et il a été mené sur les quatorze années I à XIV. Il suppose aussi que la
transcription a bien lu « germinal » et « mercredi », ce que rien ici ne permet
de contrôler. Ce qui trancherait entre les deux dates : l'adresse au dos du
billet, qui porte « Mademoiselle Sophie Germain / Paris » et aucune rue.}
Lagrange écrit son nom « Germaine ».

\subsection*{Les trois feuillets de Berlin : ils ne sont pas d'une seule lettre, ni d'une seule année}

Trois feuillets de Lagrange écrits de Berlin sont reliés ici. Ils ne sont pas
adressés à Sophie Germain --- le premier est daté du 10 juillet 1776, trois
mois après sa naissance --- mais à un correspondant de Turin que les feuillets
ne nomment pas, et qui est de Denina, du docteur Cigna, de Richeri secrétaire
du comte Fontana, et du père de Lagrange, resté à Turin.\footnote{La
transcription le note déjà : « Elle n'est pas adressée à Sophie Germain. » On
le répète parce que la reliure, elle, les présente au milieu de sa
correspondance, et qu'un catalogue qui date le volume du XVIIIe siècle le doit
sans doute à ces trois feuillets.}

Le volume les donne à la file, et la reliure invite à les prendre pour une
seule lettre. Ils sont trois, et de trois moments différents. Le contenu le
montre.

\begin{itemize}
  \item \emph{Feuillet 24}, signé « De la Grange », daté « à Berlin ce 10
    Juillet 1776 ». Le premier feuillet manque : le texte commence sans
    suscription.
  \item \emph{Feuillet 25}, sans suscription ni signature : le désagrément de
    Denina, les lois du pays, Galilée, et l'envoi d'« une medaille de
    l'Imperatrice de Russie frappée à l'occasion du dernier jubilé
    Academique ». Le jubilé de l'Académie de Saint-Pétersbourg est celui de
    1776, et l'affaire qui vaut à Denina son « désagrément » --- un ouvrage
    dont « il n'a que du chagrin », opposé à « son histoire d'Italie » qui ne
    lui a valu que de la satisfaction --- est celle de 1777. Ce feuillet est
    donc de 1777 ou peu après.
  \item \emph{Feuillet 26}, commençant « Monsieur », sans signature, le
    feuillet suivant manquant : les estampes du Roi, l'histoire diplomatique
    de Weguelin, et « les grandes armées qui sont maintenant en campagne »,
    avec cette phrase --- « comme on assure que le Roi est entré en Boheme on
    ne tardera pas à apprendre quelque chose d'interessant ». Frédéric II
    entre en Bohême en juillet 1778, au début de la guerre de Succession de
    Bavière. Ce feuillet est de l'été 1778.
\end{itemize}

Les trois feuillets couvrent donc à peu près 1776--1778, et la date portée sur
le premier ne vaut que pour lui.\footnote{Les rapprochements du dernier
paragraphe --- le jubilé de Saint-Pétersbourg, l'affaire Denina, l'entrée en
Bohême --- sont ceux de cette lecture. Ils reposent sur des faits extérieurs au
volume et vérifiables ailleurs, et non sur une pièce du volume ; c'est
pourquoi ils sont donnés comme une datation par le contenu, qu'une collation
avec la correspondance imprimée de Lagrange confirmerait ou renverserait en
une heure. Cette lecture n'a consulté aucune édition.} On notera aussi que le
millésime du premier feuillet est lui-même douteux : la transcription lit le
troisième chiffre « 7 ou 5 », et n'exclut 1756 que parce que Lagrange n'était
pas à Berlin cette année-là.

\section*{« Cet élève Gallois » : la lettre du 18 avril 1831}

\pagerange{27}{29}
Quatre pages de sa main, signées, les seules du volume. Elle meurt dix
semaines plus tard.

La lettre est d'abord celle d'une femme malade qui règle des affaires de
correspondance : des paquets qui n'arrivent pas, une lettre remise à Madame
Cauchy, Madame Renaud qui avait perdu son adresse, un domestique qui donne
reçu. Puis, sans transition :
\begin{quote}
Ma santé est dans un état affreux, une mort prompte seroit pour moi un
soulagement car je souffre des maux inconnus et qui ne me laissent pas un seul
moment de repos, je voulois lire au moins le 2me v. de la revue mais je ne
puis, je reste enfermée.
\end{quote}
Puis, dans un seul paragraphe, l'état de la science comme elle le voit :
\begin{quote}
décidément il y a un tort sur tout ce qui tient aux mathématiques votre
préoccupation celle de Cauchy, la mort de Mr fourier, pour […] cet élève
Gallois qui malgré son impertinence annonçoit des dispositions heureuses en a
tant fait qu'il a été chassé de l'école normale, il est sans fortune et sa mere
en a fort peu.
\end{quote}
Et la suite, sur la mère de Galois quittant sa maison et se plaçant dame de
compagnie, sur l'insolence dont Libri lui-même avait reçu « un échantillon
après votre lecture à l'académie », et pour finir : « On dit qu'il devient tout
à fait fou et je le crains. »

Ce que cette page établit, sans interprétation : qu'en avril 1831 elle tient le
compte de ce qui est arrivé aux mathématiques françaises en dix mois --- Libri
en exil, Cauchy parti, Fourier mort l'année précédente, Galois chassé --- et
qu'elle range le dernier avec les autres, c'est-à-dire qu'elle le tient pour une
perte pour la science et non pour un fait divers. Elle écrit son nom
« Gallois », avec deux $l$.\footnote{La transcription signale l'orthographe et
note qu'elle ne se rencontre nulle part ailleurs dans le volume, ce qui est
vrai puisque le nom n'y paraît qu'une fois. La lettre est antérieure de trois
semaines à l'arrestation de Galois, le 10 mai 1831 : la précision est de cette
lecture et n'est pas dans le feuillet.} Aucune mathématique n'est faite dans
cette lettre ; c'est pourtant la pièce du volume qu'un mathématicien
d'aujourd'hui lira le plus attentivement.

\section*{Dater ce qui n'est pas daté : les feuillets d'adresse}

\pagerange{6}{42}
Onze pièces du volume ne portent pas de date, ou ne portent qu'un jour de la
semaine. Le volume fournit lui-même de quoi les situer, et c'est son apport le
moins visible et le plus sûr : les couvertures.

Ses adresses, dans l'ordre où le volume les date :
\begin{itemize}
  \item 4 novembre 1797, Lalande : « à la Citoyenne / Sophie germain / rue Ste
    croix de la Bretonnerie n° 18 ».
  \item 14 mai 1810, Delambre : même rue, « maison de M. Doulcet d'Egligny
    maire du 7e arrondissement », n° 25 d'une lecture douteuse.
  \item 13 mars 1814, l'inconnu : même rue, numéro illisible sous deux larges
    boucles d'encre.
  \item 1820 ou 1821 --- 1827, Delambre, Fourier, Libri : « rue de Braque au
    marais », n° 41.\footnote{Delambre, en juillet 1821, hésite : « N°. 47 ou
    41 ». Les autres couvertures donnent 41, et la transcription marque le
    chiffre comme douteux sur trois d'entre elles, le 4 et le 1 étant ligaturés
    et pouvant se lire « 4 » seul.}
\end{itemize}

Il suit immédiatement que le déménagement de la rue Sainte-Croix-de-la-
Bretonnerie à la rue de Braque a eu lieu \emph{entre le 13 mars 1814 et le 21
juillet 1821}, et que toute pièce du volume dont la couverture porte « rue de
Braque » est postérieure à 1814.\footnote{La borne se resserre à septembre 1820
si l'on admet que le feuillet d'adresse 7v est la couverture de la lettre de
Fourier du 29 septembre 1820, ce que l'ordre du montage indique et ce que la
transcription affirme dans son en-tête. Cette lecture ne s'appuie que sur la
borne large, qui ne dépend d'aucune attribution de couverture : la lettre de
Delambre du 21 juillet 1821 est datée sur elle-même.} Trois des billets non
datés de Fourier sont dans ce cas.

Deux autres indices, donnés par les feuillets eux-mêmes :
\begin{itemize}
  \item \emph{Sa mère est saluée.} Deux billets de Fourier saluent « madame sa
    mère » ou l'invitent ; sa mère meurt en 1823.\footnote{La date de 1823 est
    donnée par la transcription en note, comme un fait extérieur au feuillet ;
    elle n'est écrite nulle part dans le volume. Toute la borne supérieure en
    dépend, et cette lecture la signale plutôt que de la faire passer pour une
    lecture.} Ces billets sont donc antérieurs à 1823, et postérieurs à 1814
    par leur couverture : ils tombent dans la fenêtre 1820--1823.
  \item \emph{Libri est nommé.} Deux autres billets de Fourier parlent d'un
    mémoire de Libri et d'une recommandation que Libri lui a écrite ; ils sont
    donc postérieurs à l'entrée de Libri dans le cercle parisien, et
    antérieurs à la mort de Fourier en mai 1830. La lettre de Florence du 17
    novembre 1825 rend 1825--1826 vraisemblable, sans plus.
\end{itemize}

Deux couvertures ne portent ni rue ni ville : les billets qu'elles fermaient
ont été portés à la main, et elles ne datent rien.

\section*{Ce qui n'est pas des mathématiques}

\pagerange{6}{41}

\subsection*{Les billets du centre}

La lettre de Delambre du 21 juillet 1821, après le paragraphe sur son mémoire,
consacre deux pages à une question d'accès. Elle a demandé à assister aux
séances publiques de l'Académie ; il explique pourquoi c'est difficile. Le
Bureau dispose de quarante billets et les académiciens sont soixante-quinze ;
lui-même en a « tout au plus dix » par séance et se fait une loi de les donner
aux confrères qui en demandent pour leurs femmes ; heureusement, ajoute-t-il,
« leurs femmes ne montrent pas un grand empressement en sorte que jusqu'ici je
n'ai pas eu le chagrin d'en refuser une seule ». Il promet de lui en réserver
un au mois de mars suivant.

Le document se passe de commentaire. On notera seulement qu'il est écrit par le
secrétaire perpétuel, dans la même lettre où il lui dit que personne n'est en
état de juger son mémoire.

\subsection*{Lalande, 4 novembre 1797}

Le billet le plus ancien du volume qui la concerne, et le plus étrange : des
excuses, le lendemain d'une visite.
\begin{quote}
il etoit difficile mademoiselle de me faire sentir plus que vous ne l'avez fait
hier l'indiscretion de ma visite, et l'improbation de mes hommages. […] il ne
me reste donc qu'à vous faire des excuses de mon imprudence.
\end{quote}
Il y mêle deux détails utiles : « mon ami Cousin », qui lui avait annoncé ses
talents, et le fait qu'elle a lu « le systeme du monde de la place » et ne veut
pas lire son « abregé d'astronomie ». Elle a vingt et un ans ; le \emph{Système
du monde} de Laplace a paru l'année précédente. La moitié basse du billet est
écrite très menu et plusieurs mots n'ont pas pu être lus.

\subsection*{Choron, et la lettre du 13 mars 1814}

Le billet de Choron, du 29 janvier 1818, accompagne un exemplaire de ses
observations sur l'enseignement simultané de la lecture et de l'écriture en
musique. Il n'a rien de mathématique et tout de documentaire : Choron ne figure
pas dans la notice du volume.

La lettre non signée du 13 mars 1814 est, de toutes les pièces, la plus
personnelle : deux « titres » dont l'auteur a touché le montant, Saint-Denis et
Montbrison, son cabinet contre sa bibliothèque à elle, une sœur dont elle est
« la seule ressource », et cette phrase où l'invasion de 1814 passe sans être
nommée --- « Dans Celui où nous nous trouvons il faut plus de mistère pour
respirer en liberté ». Elle est signée d'un paraphe sans lettre lisible. Son
intérêt pour cette lecture est entier et indirect : c'est sur sa couverture, et
après elle, qu'on a écrit la seule formule du volume.

\section*{Ce que les transcriptions laissent en suspens}

\pagerange{6}{42}
Ce qui suit n'est pas une liste des mots douteux --- la transcription les porte
tous --- mais des seuls points dont dépend quelque chose de ce qui précède.

\begin{itemize}
  \item \emph{Quelle couverture va avec quelle lettre.} L'attachement d'un
    feuillet d'adresse à une lettre repose sur l'ordre du montage et sur rien
    d'écrit. La transcription du premier lot l'affirme dans son en-tête pour le
    feuillet 7v --- couverture de la lettre du 29 septembre 1820 --- et son
    corps le place, en suivant l'ordre des vues, en tête de la section de la
    lettre du 30 janvier 1827. Il n'y a pas là de contradiction sur le fait,
    mais il y a deux présentations du même feuillet, et la datation du
    déménagement en dépend --- raison pour laquelle cette lecture s'appuie
    ailleurs.
  \item \emph{La main de la note mathématique.} Non tranchée, et non
    tranchable sur quinze mots. Tout ce que cette lecture avance sur la note
    est indépendant de l'attribution.
  \item \emph{Le portrait gravé.} Fourier envoie « le portrait d'un de nos
    prédécesseurs, et l'un des plus illustres fondateurs d'une science que vous
    aimez et que vous avez augmentée ». Le nom n'est pas donné et cette lecture
    ne le devine pas. L'exemplaire, s'il est resté dans ses papiers, le dirait.
  \item \emph{Les mots que la transcription donne pour douteux et dont le sens
    dépend.} « Jaux » dans la lettre de 1831, deux fois, sans doute
    une abréviation de « journaux » ; « menues » dans la lettre de
    1814 ; « confiez » et « me mande » dans le billet
    « vendredi matin », dont dépend le rattachement de ce billet à la demande
    de Libri ; « bresslieur » et « vauborel », deux noms
    d'invitées chez Fourier, qu'aucune lecture n'a arrêtés.
  \item \emph{Le numéro de la rue sur la couverture de 1814}, recouvert de deux
    boucles d'encre tracées après coup. C'est le seul document qui pourrait
    resserrer la date du déménagement par en bas.
\end{itemize}

\end{document}
